\section{单圈因子化}

现在我们可以在单圈阶构造因子化公式。在无限动量系或光锥上，部分子分布与劈裂函数中的动量分数被限制在 $[0,1]$。但在目前的情形中，准分布中的劈裂并不受该区域约束，它可以取遍 $[-\infty,\infty]$。于是，两种分布之间的联系通过下面的大 $P^z$ 极限下、精确到幂次修正的因子化定理体现：
\begin{equation}
    \tilde q(x, \mu^2, P^z)  = \int^1_{0} \frac{dy}{y} Z\left(\frac{x}{y},\frac{\mu}{P^z}\right) q(y, \mu^2) +  {\cal O}\left(\Lambda^2/(P^z)^2,  M^2/(P^z)^2\right)\ ,
\end{equation}
其中积分范围由光锥上夸克分布 $q(y)$ 的支集决定。

$Z$ 因子具有按 $\alpha_s$ 展开的微扰展开式
\begin{equation}
             Z(\xi, P^z/\mu) = \delta (\xi-1) + \frac{\alpha_s}{2\pi} Z^{(1)}(\xi, P^z/\mu) + \dots
\end{equation}
在继续之前，必须指出一个重要事实：如上面的式~(\ref{unpolvertex})与式~(\ref{unpolself})所示，准夸克分布的单圈修正中存在一个与轴向规范奇异性相耦合的线性发散，它在 $\xi=1$ 处表现为二重极点。正如上一节所述，若选择协变规范（例如 Feynman 规范），图~\ref{onelooppdf}中的各图不会产生轴向奇异性；但此时会出现额外的含规范链接的图，奇异性正来自这些图。显然，二重极点结构在维数正规化中不出现，但在本文使用的截断正规化中确实存在。与光锥部分子分布的情形不同——那里在 $\xi=1$ 处至多遇到单极点，可以用加函数约定（plus-prescription）恰当地正规化——加函数约定无法完全正规化二重极点的发散：一般而言，它只能把二重极点约化为单极点，而对 $\xi$ 积分后该单极点仍给出奇异的贡献。然而事实证明，在我们的情形中，这一奇异性可以用主值（principal value）约定消除，这对应于对 Wilson 线自能的一种正规化。虽然这个线性发散也可以被单独挑出（在格点上可以通过固定 $P^z$ 与 $x$、改变格点间距来实现）并加以减除~\cite{footnote1}，但为了简化从格点数据提取光锥分布的过程，我们提议把这个线性发散项保留在因子化公式之内。

现在我们可以写下连接准夸克分布与光锥夸克分布的匹配因子。对 $\xi>1$，有
\begin{equation}
   Z^{(1)}(\xi)/C_F = \left(\frac{1+\xi^2}{1-\xi}\right)\ln \frac{\xi}{\xi-1} + 1 + \frac{1}{(1-\xi)^2}\frac{\mu}{P^z}  \ ,
\end{equation}
而对 $0<\xi<1$
\begin{equation}
   Z^{(1)}(\xi)/C_F = \left(\frac{1+\xi^2}{1-\xi}\right)\ln \frac{(P^z)^2}{\mu^2}+\left(\frac{1+\xi^2}{1-\xi}\right)\ln\big[{4\xi(1-\xi)}\big]-\frac{2\xi}{1-\xi}+1+\frac{\mu}{(1-\xi)^2P^z} \ ,
\end{equation}
对 $\xi<0$
\begin{equation}
   Z^{(1)}(\xi)/C_F = \left(\frac{1+\xi^2}{1-\xi}\right)\ln \frac{\xi-1}{\xi}-1 +\frac{\mu}{(1-\xi)^2P^z} \ .
\end{equation}
在 $\xi=1$ 附近，还有一项来自自能修正的额外贡献
\begin{equation}
    Z^{(1)}(\xi) = \delta Z^{(1)}(2\pi/\alpha_s) \delta(\xi-1)\ ,
\end{equation}
它可以从上面的式~(\ref{unpolself})与式~(\ref{unpolselfIMF})中提取出来，并为 $\xi=1$ 处的奇异性提供加函数正规化。这样一来，$\tilde q(x,\mu^2,P^z)$ 中对 $P^z$ 的大对数依赖便可通过上述匹配条件转化为重正化标度依赖。
在格点上，还须用标准方法~\cite{ren}把匹配重新计算到常数精度。

到目前为止，我们只考虑了夸克的贡献。也可以从反夸克出发做单圈计算。此时，$\tilde q(x, \mu^2, P^z)$ 中还会有来自 $\bar q(x)$ 的贡献。然而，反夸克分布具有性质
\begin{equation}
     \bar q(x) = - q(-x)\ ,
\end{equation}
即它与负 $x$ 处的夸克分布相联系。把夸克与反夸克贡献都包括进来，便得到积分限延拓为从 $-1$ 到 $+1$ 的如下因子化公式
\begin{equation}\label{1loopfact}
    \tilde q(x, \mu^2, P^z)  = \int^1_{-1} \frac{dy}{|y|} Z\left(\frac{x}{y},\frac{\mu}{P^z}\right) q(y, \mu^2) +  {\cal O}\left(\Lambda^2/(P^z)^2,  M^2/(P^z)^2\right)\ ,
\end{equation}
其中负 $y$ 区域包含反夸克贡献。上式就是完整的单圈因子化定理，它取代了文献~\cite{Ji:2013dva}中的式 (11) 与其中的 $Z$ 因子。

我们在单圈水平上构造了连接准部分子分布与光锥部分子分布的因子化公式。当然，这样的公式是否对所有阶都存在仍有待证明。借助该因子化公式，人们可以在格点上测量不含时间的非局域夸克关联函数 $\overline{\psi}(z)\gamma^zL(z,0)\psi(0)$——其中
$L(z,z_0) = \exp\left(-ig\int^{z}_{z_0} dz' A^z(z') \right)$——以计算准部分子分布，并让核子处于越来越大的 $P^z$ 态（最大约 $\sim 1/a$，$a$ 为格点间距），从而提取部分子分布 $q(x,\mu^2)$。
